% Part: lambda-calculus
% Chapter: lambda-definability
% Section: truth-values

\documentclass[../../../include/open-logic-section]{subfiles}

\begin{document}

\olfileid{lam}{ldf}{tvr}
\olsection{قيم الصدق والعلاقات}

لنا أن نرمز قيم الصدق في حساب لامبدا الخالص على الوجه الآتي:
\begin{align*}
  \fn{true} & \ident \lambd[\OLId{latin-ordinary}{x}][\lambd[\OLId{latin-ordinary}{y}][\OLId{latin-ordinary}{x}]]\\
  \fn{false} & \ident \lambd[\OLId{latin-ordinary}{x}][\lambd[\OLId{latin-ordinary}{y}][\OLId{latin-ordinary}{y}]]
\end{align*}

ووجه هذا الترميز أن قيمة الصدق تقوم مقام \emph{دالة اختيار}: تأخذ
وسيطين ثم تردّ أحدهما. فـ~$\fn{true}$ تختار الأول، و$\fn{false}$
تختار الثاني. ولذلك يختزل~$\fn{true}\, \OLId{latin-looped}{M} \OLId{latin-looped}{N}$ أبدًا إلى~$\OLId{latin-looped}{M}$،
ويختزل~$\fn{false}\, \OLId{latin-looped}{M} \OLId{latin-looped}{N}$ أبدًا إلى~$\OLId{latin-looped}{N}$.

\begin{defn}
تكون العلاقة~$\OLId{latin-looped}{R} \subseteq \Nat^\OLId{latin-ordinary}{k}$ !!{lambda definable} إذا وجد
حد~$\OLId{latin-looped}{R}$ يحقق
\begin{align*}
  \OLId{latin-looped}{R}\, \num{\OLId{latin-ordinary}{n}_\OLMathNumeral{1}} \dots \num{\OLId{latin-ordinary}{n}_\OLId{latin-ordinary}{k}} & \bred \fn{true}
  \intertext{متى تحقق~$R(n_1, \dots, n_k)$، و}
  \OLId{latin-looped}{R}\, \num{\OLId{latin-ordinary}{n}_\OLMathNumeral{1}} \dots \num{\OLId{latin-ordinary}{n}_\OLId{latin-ordinary}{k}} & \bred \fn{false}
\end{align*}
إذا لم تتحقق.
\end{defn}

ومن أمثلة ذلك العلاقة~$\fn{IsZero} = \{\OLMathNumeral{0}\}$، وهي لا تتحقق إلا
في~$\OLMathNumeral{0}$؛ فإنها !!{lambda definable} بالحد
\[
  \fn{IsZero} \ident \lambd[\OLId{latin-ordinary}{n}][\OLId{latin-ordinary}{n} (\lambd[\OLId{latin-ordinary}{x}][\fn{false}])\, \fn{true}].
\]
وبيان عمله أن أعداد تشيرش مكررات، أي دوال تطبق وسيطها الأول~$\OLId{latin-ordinary}{n}$ مرة
في وسيطها الثاني. فنجعل مبدأ التكرار~$\fn{true}$، ونجعل كل خطوة تردّ
$\fn{false}$ بصرف النظر عما ردّته الخطوة قبلها. وتطبق هذه الخطوة في
القيمة المبدئية~$\OLId{latin-ordinary}{n}$ مرة؛ فيبقى الناتج~$\fn{true}$ إذا وفقط إذا لم
تقع خطوة أصلًا، أي عندما~$\OLId{latin-ordinary}{n} = \OLMathNumeral{0}$.

وبهذا التمثيل نفسه نضع دوال صدق أخرى. فالحدان الآتيان يمثلان النفي
والاقتران:
\begin{align*}
  \fn{Not} & \ident \lambd[\OLId{latin-ordinary}{x}][\OLId{latin-ordinary}{x}\, \fn{false}\, \fn{true}]\\
  \fn{And} & \ident \lambd[\OLId{latin-ordinary}{x}][\lambd[\OLId{latin-ordinary}{y}][xy \,\fn{false}]]
\end{align*}
تأخذ الدالة ``$\fn{Not}$'' وسيطًا واحدًا، فتردّ~$\fn{true}$ إذا
كان~$\fn{false}$، وتردّ~$\fn{false}$ إذا كان~$\fn{true}$. وأما
``$\fn{And}$'' فتأخذ قيمتي صدق، وينبغي أن تردّ~$\fn{true}$ إذا وفقط
إذا كان كل من الوسيطين~$\fn{true}$. وقد مر أن قيم الصدق دوال اختيار؛
فإذا كان~$\OLId{latin-ordinary}{x}$ قيمة صدق وطبق في وسيطين، ردّ الأول إذا كان~$\OLId{latin-ordinary}{x}$ هو
$\fn{true}$، وردّ الثاني فيما عدا ذلك. وتأخذ~$\fn{And}$ وسيطيها
$\OLId{latin-ordinary}{x}$ و$\OLId{latin-ordinary}{y}$، ثم تجعل~$\OLId{latin-ordinary}{y}$ و$\fn{false}$ وسيطي~$\OLId{latin-ordinary}{x}$. فإذا افترضنا أن
$\OLId{latin-ordinary}{x}$ قيمة صدق، كان الناتج~$\OLId{latin-ordinary}{y}$ متى كان~$\OLId{latin-ordinary}{x}$ هو~$\fn{true}$، وكان
$\fn{false}$ متى كان~$\OLId{latin-ordinary}{x}$ هو~$\fn{false}$؛ وهذا عين المطلوب.

وإنما افترضنا في هذا البيان ألا تطبق~$\fn{And}$ إلا في قيم الصدق.
فإن طبقتها في حدود أخرى، لم يلزم أن يكون الناتج، أي الصورة السوية إن
كانت، قيمة صدق.

\begin{prob}
  عرّف الدالتين~$\fn{Or}$ و$\fn{Xor}$ اللتين تمثلان دالتي صدق الفصل
  الشامل والفصل الحصري، مستعملًا ترميز قيم الصدق بحدود~$\lambd$.
\end{prob}

\end{document}
